% GPT-6 Astra Pro; effort not exposed; /root/replacement; 22 November 2026.
\section[HD-F1-R27: the common character tensor is nonzero]{HD-F1-R27: effective target surfaces exclude the first rank-two character alternative}
\label{sec:followup-rm-rank-two-surface-positivity}

For every remaining actual adjacent object of R21--R26
at ranks $(r_A,r_T)=(3,2)$, the tensor $M$ in the
unique decomposition (F1R26.14) is nonzero. Thus the
first alternative of (F1R26.21) is excluded, and the
whole characters necessarily have the form
\[
 \operatorname{ch}(T)=a_0v_q+c\,p\overline p,\qquad
 \operatorname{ch}(A)=(a_0+1)v_q+c\,p\overline p
 \tag{F1R27.1}
\]
with the exact rationality, nonzero projection and
two conjugate Hodge conditions of (F1R26.16).
This conclusion holds for the arbitrary ample
polarization $H$ with respect to which the two
actual endpoints are semistable of equal slope.
It does not require $H$ to be compatible with the
real-multiplication action.

The proof uses the actual target surface defects
already forced by R22 and R25. No cohomological
factor in (F1R27.1) is asserted to be a sheaf
character, and no sheaf or coupling is constructed
by this exclusion.

\paragraph{1. The actual cycle-to-form map and a strictly positive mixed pairing.}
Write the two complex RM tangent eigenspaces as
$T_0X=T_1\oplus T_2$, each of dimension two.
Let $\omega_i$ be any invariant positive Hermitian
$(1,1)$ form on $T_i$, extended by zero on the
other block. They need not be rational forms.
For every nonzero effective surface cycle
$Z=\sum n_jZ_j$ on $X$, one has
\[
                       \int_Z\omega_1\omega_2>0.
 \tag{F1R27.2}
\]
Here and below integration on a singular component
means pullback to its smooth projective resolution.

First, $\omega_1\omega_2$ is a positive sum of
the squares of all mixed holomorphic two-forms
in $T_1^*\otimes T_2^*$, after choosing unitary
bases in the two blocks. Its pullback to a
surface is nonnegative. If the integral were
zero, every mixed holomorphic two-form would
pull back to zero on every component.
At a smooth point of a component its Pl\"ucker
bivector would therefore have the form
$u_1+u_2$, $u_i\in\bigwedge^2T_i$.
Decomposability gives $u_1u_2=0$, so its tangent
plane is $T_1$ or $T_2$.
The Gauss map on an irreducible component is
then constant with one of these values.

For completeness the consequence is an actual
map, as in the supplied DEG\_L S-22 proof.
Translate a resolution $f:Y\to Z_j\subset X$
so that a chosen point maps to zero. Every
invariant one-form annihilating the constant
tangent plane pulls back to zero. The actual
induced homomorphism
$\operatorname{Alb}(Y)\to X$ consequently has
tangent image contained in $T_i$.
Its image is positive dimensional because $f$
is nonconstant. Such an image would give a
nonzero rational homology vector in the
underlying real RM eigenspace. Choosing
$u\in F\setminus\mathbf Q$ gives on that
vector the equation $u_*z=\sigma_i(u)z$.
The matrix of $u_*$ and the vector $z$ are
rational, whereas $\sigma_i(u)$ is irrational;
any nonzero coordinate contradicts this
equation. This proves (F1R27.2) for every
choice of the positive block forms.

We also need the pointwise positivity of the
particular invariant representative of $[Z]$.
The actual comparison is Haar averaging of the
integration current:
\[
               [Z]\longmapsto
                 \overline{[Z]}:=\int_X t_x^*[Z]\,dx.
 \tag{F1R27.3}
\]
It is a closed positive translation-invariant
current. Translation invariance makes its
coefficients constant, so it is a smooth
invariant form. Translations are isotopic to
the identity, and averaging therefore preserves
its cohomology class. The exterior algebra
of invariant one-forms identifies invariant
forms with the cohomology of the real torus;
in particular, there is exactly one invariant
representative of each cohomology class.
Thus (F1R27.3) is that representative.
Equivalently, for every invariant complementary
form $\eta$, its pairing is exactly
$\int_Z\eta$. The dual positive-cone calculation
in Debarre--Ein--Lazarsfeld--Voisin, Lemma~1.5,
shows that this representative is strongly
positive, and hence weakly positive.
In particular it evaluates nonnegatively on
every complex mixed plane
$\mathbf Cu\oplus\mathbf Cv$, $u\in T_1$,
$v\in T_2$. This is a pointwise assertion on
the actual representative (F1R27.3), rather
than only one intersection-number inequality.

\paragraph{2. A two-line surface obstruction for arbitrary $H$.}
Let $\ell,m$ be real invariant $(1,1)$ forms
with
\[
       \ell H^3=mH^3=0,\qquad
       \ell+m\in H^{1,1}(T_1)\oplus H^{1,1}(T_2).
 \tag{F1R27.4}
\]
For any real $u,v$, the class
\[
              \gamma=\tfrac12(\ell^2+m^2)-uA_1-vA_2
 \tag{F1R27.5}
\]
cannot be the class of a nonzero effective
surface cycle. We prove this including all
off-block components of the arbitrary ample
form $H$.

Represent $H$ by its positive Hermitian matrix.
The zero-degree conditions in (F1R27.4) are
exactly $\operatorname{tr}(H^{-1}\ell)=
\operatorname{tr}(H^{-1}m)=0$.
Choose a block diagonal change of tangent
coordinates that makes the two diagonal
blocks of $H^{-1}$ equal to the identity.
Explicitly, for its original diagonal blocks
$K_{11},K_{22}>0$, use the change matrix
$\operatorname{diag}(K_{11}^{1/2},K_{22}^{1/2})$.
The inverse metric and the two divisor matrices
then have the form
\[
 H^{-1}=\begin{pmatrix}I&R\\R^*&I\end{pmatrix},
 \qquad \|R\|_{\rm op}<1,\qquad
 \mathsf L=\begin{pmatrix}A&Z\\Z^*&B\end{pmatrix},
 \qquad
 \mathsf M=\begin{pmatrix}A'&-Z\\-Z^*&B'\end{pmatrix}.
 \tag{F1R27.6}
\]
The strict norm bound is the Schur-complement
identity $I-R^*R>0$ for this positive matrix.
The opposite off-block entries follow from
the second assertion of (F1R27.4); no separate
block primitivity is assumed.
Put
\[
 P=A+A',\quad Q=B+B',\quad D=A-A',\quad J=B-B'.
\]
The actual two trace equations give
\[
       \operatorname{tr}(P+Q)=0,\qquad
       \operatorname{tr}(D+J)
                    =-4\operatorname{Re}\operatorname{tr}(R^*Z).
 \tag{F1R27.7}
\]

Suppose (F1R27.5) represents an effective cycle.
Its invariant representative is precisely
the displayed form, by paragraph~1. On a mixed
plane both $A_i$ vanish. The coefficient of
$\ell^2/2$ is the determinant of the restricted
two-by-two Hermitian matrix, and likewise for
$m$. Weak positivity therefore says, for all
unit vectors $a,b$ in the two blocks,
\[
   (a^*Pa)(b^*Qb)+(a^*Da)(b^*Jb)
                              \ge4|a^*Zb|^2.
 \tag{F1R27.8}
\]
This follows by adding the two determinants
and using
$AB+A'B'=(PQ+DJ)/2$ for their scalar diagonal
entries.

We next prove $Z=0$. Suppose instead $Z\ne0$.
A unitary polar decomposition, with a unitary
extension on the kernel, identifies the two
blocks so that $Z=S\ge0$ and
$\operatorname{tr}S=\|Z\|_*>0$, where the last
quantity is the sum of its singular values.
This independent unitary change on the two
blocks preserves (F1R27.7) and $\|R\|_{\rm op}$.
The two-by-two Hermitian matrix $P+Q$ has
trace zero. Choose an orthonormal basis
$e_1,e_2$ with
\[
          (P+Q)_{11}=(P+Q)_{22}=0,
                   \qquad s_i=S_{ii}>0\quad(i=1,2).
 \tag{F1R27.9}
\]
Here is the full existence argument. For a
nonzero trace-zero $P+Q$, choose its two
orthonormal eigenvectors and take their
equal-length sums and differences with an
arbitrary unit-modulus relative phase.
Both expectations of $P+Q$ are zero.
The phase can avoid the kernel line of
the nonzero positive matrix $S$, making both
expectations of $S$ positive.
When $P+Q=0$, any generic orthonormal basis
has the required property. These are changes
of tangent coordinates, not changes of the
original sheaves, action or polarization.

Write $p_i=P_{ii}$, $d_i=D_{ii}$, $j_i=J_{ii}$,
so that $Q_{ii}=-p_i$.
Equation (F1R27.8) on the diagonal mixed planes
gives
\[
                         d_i j_i\ge p_i^2+4s_i^2>0.
 \tag{F1R27.10}
\]
Thus $d_i,j_i$ have the same nonzero sign
for each $i$. Their signs cannot be opposite
for $i=1$ and $i=2$. Indeed the two off-diagonal
mixed planes would give two negative numbers
\[
                    d_1j_2\ge p_1p_2,\qquad
                    d_2j_1\ge p_1p_2.
 \tag{F1R27.11}
\]
They force $p_1p_2<0$, and their product is
at most $(p_1p_2)^2$. But the same product
equals $(d_1j_1)(d_2j_2)$, which is strictly
larger than $(p_1p_2)^2$ by (F1R27.10).
This is impossible. All four signs therefore
agree, and
\[
 \begin{split}
 |\operatorname{tr}(D+J)|
    &=\sum_{i=1}^2|d_i+j_i|
      \ge2\sum_{i=1}^2\sqrt{d_i j_i}
      \ge4(s_1+s_2)=4\|Z\|_*.
 \end{split}
 \tag{F1R27.12}
\]
On the other hand, (F1R27.7) and the elementary
singular-value estimate give
\[
 |\operatorname{tr}(D+J)|
   \le4\|R\|_{\rm op}\|Z\|_*<4\|Z\|_*.
 \tag{F1R27.13}
\]
For that estimate, write a singular-value
decomposition of $Z$ and bound each resulting
diagonal matrix coefficient of $R$ by its
operator norm. This explicitly proves the
inequality being used. The contradiction
proves $Z=0$.

Now the two individual zero-slope equations
in (F1R27.6) reduce to
$\operatorname{tr}A+\operatorname{tr}B=0$
and
$\operatorname{tr}A'+\operatorname{tr}B'=0$.
Let $\omega_i$ be the standard positive
block forms in these coordinates, and let
$\Omega$ be the positive product volume form.
Direct exterior multiplication gives
\[
 \begin{split}
 \ell^2\omega_1\omega_2
       &=2\,\operatorname{tr}A\operatorname{tr}B\,\Omega,\\
 m^2\omega_1\omega_2
       &=2\,\operatorname{tr}A'\operatorname{tr}B'\,\Omega,\\
 A_i\omega_1\omega_2&=0.
 \end{split}
 \tag{F1R27.14}
\]
It follows that
\[
       \gamma\omega_1\omega_2
          =-\bigl((\operatorname{tr}A)^2+
                    (\operatorname{tr}A')^2\bigr)\Omega\le0.
 \tag{F1R27.15}
\]
This contradicts (F1R27.2) for a nonzero
effective cycle, and proves the obstruction
(F1R27.4)--(F1R27.5).

\paragraph{3. The first rank-two character has the excluded degree-two part.}
In the first alternative of R26 write
\[
       \operatorname{ch}(T)=a\,v_1v_2+bv_2+v_1\overline b,
                    \qquad b\in V_1\cap S_1^0.
 \tag{F1R27.16}
\]
The notation here is exactly the Mukai-space
notation of R26. Decompose
$b=b_0+\beta+b_2A_1$, with
$\beta\in H^{1,1}(T_1)$.
Its orthogonality to $v_1=1-qA_1$ says
$b_2-qb_0=0$. Thus
\[
             b=b_0(1+qA_1)+\beta.
 \tag{F1R27.17}
\]
The conjugate factor has the corresponding
expression in the second block. Multiplying
the actual cohomology classes in (F1R27.16)
now gives
\[
 \begin{split}
 c_1(T)&=\beta+\overline\beta
                  \in H^{1,1}(T_1)\oplus H^{1,1}(T_2),\\
 \operatorname{ch}_2(T)
     &=q(-a+b_0-\overline b_0)A_1
           +q(-a-b_0+\overline b_0)A_2.
 \end{split}
 \tag{F1R27.18}
\]
The bars mean the scalar-extension map of R26,
with its separately imposed Hodge conditions.
The resulting whole character is rational
and real, so the displayed degree-two
coefficients are real. In particular this
degree-two term has no mixed-block component.
This calculation uses the complete form
(F1R27.16), not an assumed truncation.

\paragraph{4. The same actual sheaves supply a nonzero effective surface.}
In the common-image branch, retain the exact
sequence and the actual reflexive-hull maps
from R23:
\[
 0\longrightarrow I\longrightarrow T\longrightarrow Q\longrightarrow0,
 \qquad I\longrightarrow L=I^{**},\qquad
 Q\longrightarrow M=Q^{**}.
 \tag{F1R27.19}
\]
Both rank-one TF sheaves and both line hulls
have slope zero for the given $H$.
Put $\ell=c_1(L)$, $m=c_1(M)$.
The two actual coherent quotients have
codimension at least two, and their effective
leading surface cycles give
\[
 \begin{split}
 \mathcal D&=[L/I]_2+[M/Q]_2\ne0,\\
 [\mathcal D]&=\tfrac12(\ell^2+m^2)
                         -\operatorname{ch}_2(T),\\
 \ell+m&=c_1(T).
 \end{split}
 \tag{F1R27.20}
\]
The nonvanishing is precisely R25, proved
there through the actual pushout and line
extension cup map. The character identity
here follows from the displayed short exact
sequences in coherent $K_0$; the leading
Chern character of each quotient is its
support cycle with generic lengths.
Equations (F1R27.18)--(F1R27.20) satisfy
every hypothesis of paragraph~2, with this
actual nonzero effective cycle. They are
therefore contradictory.

In the common-quotient branch, R22 supplies
the actual equivariant isomorphism
$T^{**}\simeq L^{\oplus2}$ and the exact
sequence
\[
              0\longrightarrow T\longrightarrow
                    L^{\oplus2}\longrightarrow Z_T\longrightarrow0,
          \qquad \mu_H(L)=0,\qquad [Z_T]_2\ne0.
 \tag{F1R27.21}
\]
Its nonzero surface part follows from the
actual degree-two predecessor in (F1R22.15)
and the proved cancellation rank at least
three. The same $K_0$ maps give
\[
       [Z_T]_2=\ell^2-\operatorname{ch}_2(T),
                    \qquad 2\ell=c_1(T).
 \tag{F1R27.22}
\]
Apply paragraph~2 with the two real forms
$\ell,m=\ell$. Again (F1R27.18) provides
the block condition and the pure degree-two
term, while (F1R27.21) provides the actual
nonzero effective cycle. This is the same
contradiction.

Both alternatives in the actual rank-one
factorization of R21 have now been treated.
Thus $M=0$ in (F1R26.14) is impossible for
these endpoints, proving (F1R27.1).

\paragraph{Framework, source scope and what is retained.}
The determinant line twist and the coherent
point-action conjugation are exactly
(F1R21.2)--(F1R21.3); no slope condition is
changed by a coordinate choice. The maps
from actual sheaves to the two surface
cycles are (F1R27.19) and (F1R27.21),
with their coherent $K_0$ and leading
cycle-class maps. The average (F1R27.3)
identifies their particular cohomological
representatives and supplies every mixed-plane
inequality needed above.
The published positive-cone input is
Debarre--Ein--Lazarsfeld--Voisin,
\emph{Compositio Mathematica} 147 (2011),
1793--1818, Definition~1.1, Remark~1.3 and
the complete Lemma~1.5 proof.
The strictly positive RM pairing is proved
in paragraph~1 by the actual Albanese map,
using the complete supplied DEG\_L S-22
tangent-space argument.
The two-by-two matrix argument is proved
in full for arbitrary $H$ in paragraph~2.

Using only a single mixed intersection
number before removing the off-block
divisor entries would be insufficient
for an arbitrary $H$. The pointwise
inequalities (F1R27.8), and specifically
their off-diagonal sign comparison, are
what proves the vanishing of those entries.
No unjustified replacement of $H$ by a
block polarization is made.
The surviving case in (F1R27.1) retains
the actual nonzero tensor, both target
surface defects and the original connecting
class $\delta\in\operatorname{Ext}^2_G(T,A)$.
It is a strict exclusion of the other
whole-character alternative, not an
existence claim for the one that remains.
